% =====================================================================
% Motivation -- LaTeX copy of research/phd/Motivation.md
% Converted 2026-10-31 from the vault version of that day. No wording changed.
% Citations: \cite{} with your Zotero citekeys, plus one placeholder key
% (Chatterji2025, not in Zotero). All entries are in state-of-the-art.bib.
% =====================================================================

\section{Motivation}

Artificial intelligence is undergoing a fundamental transformation in its role within human societies~\cite{rahwanMachineBehaviour2019}. AI systems are currently being developed with increasing autonomy and capabilities, leaning towards the edge of operational independence, and, further still, approaching full agency~\cite{hanSocialPhysicsAge2026}. The resulting increase in autonomy and trust leads humans to delegate increasingly open-ended tasks to AI agents, which progressively take initiative in engaging, adapting, and influencing outcomes alongside humans, rather than merely executing predefined instructions~\cite{hanSocialPhysicsAge2026}. This drastic shift in AI's role fundamentally changes human-AI relations, not just at the dyadic level, but also restructures the social dynamics of human populations.

AI systems are becoming deeply embedded in human societies, spanning economic and institutional processes to everyday life~\cite{rahwanMachineBehaviour2019}. The recent large language model revolution marks a turning point in this trajectory, enabling the mass deployment of AI capable of complex, context-dependent interactions with humans: by July 2025, ChatGPT alone had been adopted by around 10\% of the world's adult population~\cite{Chatterji2025}. This progressive embedding is slowly reshaping the structure of societies, transitioning from purely human-centered systems to hybrid environments where humans and AI co-exist and collaborate within shared systems~\cite{hanSocialPhysicsAge2026}. At this scale, human-AI interactions are no longer episodic but rather systematic, with millions of simultaneous interactions shaping daily life~\cite{pedreschiHumanAICoevolution2025}. Alongside this, AI systems are increasingly interacting with one another, introducing a new layer of artificial-to-artificial exchanges~\cite{dafoeOpenProblemsCooperative2020}. Together, these interactions form complex networks of humans and artificial actors from which collective patterns emerge. These patterns reshape social dynamics in ways that are difficult to predict or control.

In order to properly function cohesively, hybrid societies must rely on cooperation between humans and AI agents~\cite{dafoeCooperativeAIMachines2021}. This need for cooperation is not merely functional, but rather a fundamental prerequisite for the stable co-existence of agents with potentially diverging interests, a problem distinct from aligning an AI system with the person it serves~\cite{dafoeOpenProblemsCooperative2020}. Yet cooperation is far from trivial: agents compete for resources, creating an inherent tension between self-interest and collective benefit. This raises a fundamental question: how can cooperation emerge and sustain itself when agents with competing interests share the same environment? Artificial agents already change how much humans cooperate, even when they follow fixed rules, and whether they raise or lower it depends on the game played, on their share of the population and on the risk involved~\cite{guoFacilitatingCooperationHumanagent2023, fuOptimalIntegrationIntelligent2026, terruchaArtCompensationHow2024}.

Answering this requires moving beyond population-level observations and examining how individual agents make decisions, specifically understanding their reasoning processes. Whether an artificial agent helps, and which design helps most, already depends on how the humans around it learn~\cite{sharmaSmallBotsBig2023, zimmaroEmergenceCooperationOneshot2024}. While human reasoning is molded by cognitive biases and bounded rationality, AI reasoning is determined by its design, namely the principles and learning processes that govern how it perceives, processes, and responds to its environment. The cognitive profile of each agent, whether human or artificial, is therefore one of the factors that determine how cooperation evolves and shapes the collective dynamics of hybrid societies. Yet the AI agents now being deployed can reason about others~\cite{oecdIntroducingOECDAI2025}: language-model agents coordinate better with humans when prompted to reason about the other player~\cite{akataPlayingRepeatedGames2025}. Population models of hybrid societies, by contrast, give artificial agents a fixed policy or, at most, a simple learning rule, so their conclusions describe agents unlike those being deployed. Moreover, neither kind of reasoning stays fixed: exposure to AI changes how people think and interact with one another~\cite{riedlCognitiveSpilloverHuman2026}, and AI systems are retrained on the data their users generate~\cite{pedreschiHumanAICoevolution2025}.

While hybrid human-AI societies are already a growing reality, their full social and philosophical implications are only beginning to unfold. As AI becomes progressively embedded across societal systems, the boundaries between human and artificial agency now begin to blur~\cite{guzmanOntologicalBoundariesHumans2020}. Such a shift carries major implications, not only for how humans perceive and relate to AI at the dyadic level, but for the broader social dynamics of society as a whole. At a deeper level, this shift raises fundamental ontological questions and challenges metaphysical assumptions about the very nature of society and the mechanisms that sustain it. This underscores the urgency of fully understanding the social impact of such human-AI hybrid structures, not just to anticipate and mitigate potential conflicts, but also to guide the responsible integration of AI into human societies and support cooperation between human and artificial agents. This work investigates how the cognitive profiles of AI agents shape the emergence, stability and evolution of cooperation in hybrid human-AI societies, and whether these effects hold once the cognitive profiles of humans, and then of AI agents, change over time. This gives those who design, deploy and regulate AI agents a basis for anticipating their effects on human societies.
